\section{Mathematical Definition and Properties}
\label{sec:definition}

\subsection{Definition}

For a $k$-district congressional plan with district-level Democratic vote shares
$v_1, v_2, \ldots, v_k$, the Mean-Median Difference is:
\[
\text{MM} = \bar{v}_D - \tilde{v}_D,
\]
where $\bar{v}_D = k^{-1}\sum_{i=1}^k v_i$ is the mean and $\tilde{v}_D$ is
the median (the value $v_{(k/2)}$ for even $k$, or $v_{((k+1)/2)}$ for odd $k$,
in the sorted sequence $v_{(1)} \leq v_{(2)} \leq \cdots \leq v_{(k)}$).

Positive MM ($\bar{v}_D > \tilde{v}_D$) indicates a right-skewed distribution
of Democratic vote shares: a few districts with very high Democratic vote share
(packing) pull the mean above the median. Negative MM indicates a left-skewed
distribution (a few districts with very low Democratic vote share, consistent
with packing Democratic voters at very low margins in Republican-advantaged districts).

\subsection{Range and Interpretation}

The range of MM is approximately $[-0.10, +0.10]$ for congressional delegations;
values outside this range are rare and typically indicate either very small $k$
or severe packing. For NC ($k=14$), one district moving from 55\% to 80\%
Democratic vote share shifts MM by approximately $0.0179 / k \approx 0.013$,
setting the practical resolution.

The mean-median relationship from statistical distribution theory provides
interpretive guidance: for a symmetric distribution, $\text{MM} = 0$; positive
skewness implies positive MM; negative skewness implies negative MM. Because
geographic sorting of Democratic voters into urban cores produces a right-skewed
distribution of Democratic vote shares (a few very high-margin urban districts,
many moderate-to-low-margin suburban and rural districts), baseline MM for any
state with significant urban concentration will be positive---independently of
whether the map was gerrymandered.

\subsection{Advantages Over EG}

MM has two advantages over EG that make it preferable for single-election analysis:

\textbf{No swing model required.} MM is computed directly from observed vote shares
without applying a uniform swing transformation. EG requires either directly computing
from observed vote counts (a function of who actually won) or applying a swing model
to estimate hypothetical outcomes. MM avoids this model dependency entirely.

\textbf{Lower election-cycle sensitivity.} Because MM depends on the ranked distribution
of vote shares (how concentrated are high-Democratic-margin districts?), it is less
sensitive to uniform partisan swings than EG. A uniform 5-point Democratic surge shifts
all district vote shares by the same amount, preserving the ranked order and therefore
the mean-median gap. EG changes more dramatically because the swing shifts which districts
are won and lost, creating discontinuities in the wasted-vote calculation.

\textbf{No uncontested-district distortion.} EG calculations require special treatment
of uncontested districts (where one party receives no votes). MM treats all districts
identically regardless of contestedness, though districts with implausibly extreme
vote shares (e.g., 95\% because no opposing candidate ran) should be excluded or
corrected using standard imputation methods.

\subsection{Relationship to EG}

For a symmetric distribution of district vote shares, $\text{MM} \approx \text{EG}/(2k)$,
reflecting that both metrics increase with the skewness of the vote-share distribution.
The formal relationship depends on the shape of the distribution; empirically, for NC and
WI, $\text{MM} \approx 0.5 \times \text{EG}$ when both are computed from the same data,
consistent with the theoretical approximation for $k \approx 10$.
